\documentclass[11pt,a4paper]{article}
\usepackage[margin=24mm]{geometry}
\usepackage{amsmath,amssymb,booktabs,microtype}
\usepackage[hidelinks]{hyperref}
\newcommand{\dd}{\mathrm d}
\title{Route G2-S/D: Finite Readout and an Energy-Accounted Drive}
\author{A conditional scattering prediction, not a particle-mass fit}
\date{23 September 2026}
\begin{document}
\maketitle
\begin{abstract}
We broaden the G2-R recoil lines using a declared classical beam and
instrumental response. Rate weighting precedes conditioning; optimal
sign-classification error measures distinguishability. Separately, we
activate a homogeneous periodic modulation of the local quartic.
Its energy sidebands carry an explicit pump-work ledger. The drive
does not change the free mass tower, and can worsen discrimination.
Neither calculation is an absolute finite-encounter probability.
\end{abstract}

\section{Fixed physics and new apparatus assumptions}
Retain the G2 fields on $\mathbb M^{3,1}\times S^1$, with
\begin{equation}
 m_a^2=M_5^2+(a+\alpha)^2/R^2,\qquad
 \mu_b^2=M_D^2+b^2/R^2,\qquad j=n+l-m.
\end{equation}
The engineering card is $R=1$, $M_5=.5$, $\alpha=.25$, $M_D=5$,
$g=.05$, and $n=1$. The finite internal detector preparation remains
$w_l\propto\exp[-l^2/(2(1.2)^2)]$ for $-4\leq l\leq4$,
normalized to sum to one. All energies and momenta use the same
arbitrary mass unit as G2. Nothing here is an experimental calibration.

For resolution, condition on $m=0$ and detector mass label $r=|j|=1$.
Only $(l,j)=(0,+1),(-2,-1)$ contribute. We assume correct mass labels
and COM-frame identification, with equal unit acceptance.
Mass confusion, angular/frame uncertainty, backgrounds and efficiency
are not included. Incoming momenta are back-to-back with magnitude
\begin{equation}\label{eq:beam}
 f_{\sigma_i}(p)=
 \frac{\mathbf1_{p\geq0}e^{-(p-p_0)^2/(2\sigma_i^2)}}
 {\sqrt{2\pi}\sigma_i\Phi_{\!N}(p_0/\sigma_i)},\qquad p_0=.2 .
\end{equation}
Here $\Phi_{\!N}$ is the normal CDF; $f_0=\delta(p-p_0)$.
$p_0,\sigma_i$ are parent-Gaussian location and scale, not the
truncated distribution's actual mean and deviation.
For $a=p_0/\sigma_i$ and $\lambda=\varphi(a)/\Phi_{\!N}(a)$,
\begin{equation}
 \langle p\rangle=p_0+\sigma_i\lambda,\qquad
 \operatorname{Var}(p)=\sigma_i^2(1-a\lambda-\lambda^2).
\end{equation}
This is an \emph{incoherent classical beam ensemble}, not a coherent
four-dimensional wavepacket, duration or luminosity specification.
The common external-overlap assumption is retained.

The observed momentum estimator $x$ has response
\begin{equation}\label{eq:response}
 {\cal R}_{\sigma_d}(x\mid p_f)=
 \frac{e^{-(x-p_f)^2/(2\sigma_d^2)}}{\sqrt{2\pi}\sigma_d},
 \qquad x\in\mathbb R .
\end{equation}
A negative fitted estimator is allowed; true $p_f$ is nonnegative.
Clipping would require a different normalized response.
Scan $\sigma_i=0,.02,.1,.2,.4$ and $\sigma_d=.02,.05,.1,.2,.4$.
These are engineering choices, not fitted detector specifications.

\section{Forward model and discrimination}
Set $E_n=\sqrt{m_n^2+p^2}$, $E_l=\sqrt{\mu_l^2+p^2}$,
$E_i=E_n+E_l$. For an outgoing total energy $E$, define
\begin{equation}\label{eq:pf}
 p_f(E;m,j)=
 \frac{\sqrt{[E^2-(m_m+\mu_j)^2][E^2-(m_m-\mu_j)^2]}}{2E}
\end{equation}
only for $E>m_m+\mu_j$; closed channels have zero rate.
Undriven $E=E_i$, and the G2 rate is
\begin{equation}\label{eq:K}
 K_{lm}(p)=\frac{g^2p_f}{16\pi E_nE_lE_i}.
\end{equation}
The unnormalized readout densities are
\begin{align}
 h_j(x)&=w_{j+m-n}\int_0^\infty \dd p\,
 f_{\sigma_i}(p)K_{j+m-n,m}(p)
 {\cal R}_{\sigma_d}(x\mid p_f),\\
 Z_j&=\int_{\mathbb R}h_j(x)\dd x,\qquad Z=Z_++Z_-,\\
 P(j,x\mid m=0,r=1,\text{scattering})&=h_j(x)/Z,\qquad
 \pi_j=Z_j/Z .
\end{align}
Rate weighting must precede normalization: averaging normalized
fixed-$p$ probabilities first generally gives a different result.
Where the denominator is positive,
$P(j\mid x,m,r)=h_j(x)/(h_+(x)+h_-(x))$.

For equal classification costs, deciding the larger $h_j(x)$ minimizes
error pointwise. Thus
\begin{equation}\label{eq:bayes}
 \boxed{e_*=\frac1Z\int_{\mathbb R}\min(h_+,h_-)\,\dd x
 =\frac12\left[1-\frac{\|h_+-h_-\|_1}{Z}\right].}
\end{equation}
The no-momentum baseline is $e_{\rm prior}=\min(\pi_+,\pi_-)$,
not necessarily $1/2$. We call $e_*\leq .05$ \emph{resolved} solely
as an engineering decision criterion. It is not a discovery
significance or universal two-peak criterion. Unknown parameters need
nuisance propagation; testing a theory also needs finite event counts.

For $\sigma_i=0$, the undriven result is an exact two-Gaussian
mixture: $u_+=1.02175457609$, $u_-=1.36191317264$.
With $\Delta=u_--u_+>0$, its Bayes boundary and error are
\begin{align}
 x_*&=\frac{u_++u_-}{2}
       +\frac{\sigma_d^2}{\Delta}\log\frac{\pi_+}{\pi_-},\\
 e_*&=\pi_+\left[1-\Phi_{\!N}\left(\frac{x_*-u_+}{\sigma_d}\right)\right]
       +\pi_-\Phi_{\!N}\left(\frac{x_*-u_-}{\sigma_d}\right).
\end{align}
This supplies an independent analytic regression.
Differentiating $E_i(p)=E_f(p_f)$ gives
\begin{equation}\label{eq:slope}
 \frac{\dd p_f}{\dd p}
 =\frac{p(1/E_n+1/E_l)}{p_f(1/E'_m+1/E'_j)},\qquad
 \sigma_{f,j}^2\simeq\sigma_d^2+
 \left(\frac{\dd p_f}{\dd p}\right)^2\sigma_i^2 .
\end{equation}
The width approximation applies only to a narrow beam far from $p=0$
and thresholds. Our computation uses the nonlinear, rate-weighted
integral, not a width fit. Extra Gaussian detector noise cannot improve
$e_*$: convolution contracts the $L^1$ norm in \eqref{eq:bayes}.
Changing $\sigma_i$ changes selected priors and kinematics, so no
analogous monotonicity claim is made.

\section{Activating a powered local time-dependent background}
The new branch replaces only the interaction:
\begin{equation}\label{eq:drive}
 S_{\rm int}=-\int\dd^4x\,R\,\dd\theta\,
 \kappa_5[1+\epsilon\cos(\Omega t)]|\Phi|^2|X|^2,\qquad
 g(t)=g[1+\epsilon\cos(\Omega t)] .
\end{equation}
This is a new \emph{prescribed external pump}, not a deduction from
the motivating paper or a stabilized vacuum.
Choose $\epsilon=.5$, $\Omega=.2$, and inspect $\Omega=.1,.4$
without optimizing a desired line. For $0\leq\epsilon<1$ the
classical potential is nonnegative at all times. The quadratic action
and free mass eigenvalues remain unchanged. This does not establish
loop stability or a five-dimensional UV completion.

Spatial and circle homogeneity preserve their momenta and both global
charges. The pump rest frame is the incoming COM frame, a physical
preferred frame, not Lorentz-invariant energy injection. Expand
\begin{equation}
 g(t)=\sum_qg_qe^{-\mathrm iq\Omega t},\qquad
 g_0=g,\quad g_{\pm1}=g\epsilon/2 .
\end{equation}
At first order,
\begin{align}
 S_{fi}^{(1)}&\ \propto\
 -\mathrm i(2\pi)^3\delta^3(\boldsymbol P_f-\boldsymbol P_i)
 \delta_{n+l,m+j}
 \sum_{q=0,\pm1}g_q\,2\pi\delta(E_f-E_i-q\Omega),\\
 E_{lmq}&=E_n+E_l+q\Omega,\qquad
 K_{lmq}(p)=
 \boxed{\frac{|g_q|^2p_f(E_{lmq};m,j)}
 {16\pi E_nE_lE_{lmq}}}.
 \label{eq:drivenK}
\end{align}
Our $q>0$ means energy absorbed from the pump. The rate follows from
initial normalization $(4E_nE_l)^{-1}$ and integrated two-body phase
space $p_f/(4\pi E_{lmq})$. Do not replace the initial energies by
outgoing energies. For $q=0$ it is exactly \eqref{eq:K}.
Reference \cite{floquet} supplies the established sideband framework;
the specific relativistic rate is derived here using \cite{pdg}'s
state/phase-space convention.

Only $q=0,\pm1$ occur at this Born order; more vertices can generate
higher sidebands. This is a long-time weak-scattering rate, not a
unitarized all-orders Floquet solution. A finite rectangular time
window instead gives the energy amplitude
\begin{equation}
 \sum_q g_q F_T(E_f-E_i-q\Omega),\qquad
 F_T(\omega)=\frac{2\sin(\omega T/2)}{\omega},\quad F_T(0)=T .
\end{equation}
Finite-$T$ cross terms cannot be discarded in general.
The rate limit uses $|F_T(\omega)|^2/T\to2\pi\delta(\omega)$
distributionally for distinct fixed $\Omega>0$. The
$\Omega\to0$ and $T\to\infty$ limits do not commute:
at exactly zero frequency the amplitude is $g(1+\epsilon)$, not
three incoherent terms. Coherent beams or finite pulses need their
actual amplitude-level envelopes.

\subsection*{Work ledger and event labels}
For $H(t)=H_0+g(t)V$, unitary evolution implies
\begin{equation}
 \frac{\dd}{\dd t}\langle H(t)\rangle
 =\left\langle\frac{\partial H}{\partial t}\right\rangle
 =\dot g(t)\langle V\rangle .
\end{equation}
Matter is not energetically closed: each Born sideband has
$\Delta E_{\rm matter}=q\Omega=-\Delta E_{\rm pump}$.
With $W_{lmq}=w_lK_{lmq}$,
\begin{equation}
 Z_D=\sum_{lmq}W_{lmq},\quad
 {\cal K}_{\rm work}=\sum_{lmq}W_{lmq}q\Omega,\quad
 \langle W_{\rm pump}\rangle_{\rm scattering}
 ={\cal K}_{\rm work}/Z_D .
\end{equation}
The work coefficient is not a power; luminosity/overlap is missing.
No quantum pump apparatus, depletion or backreaction is derived.
$m=n,q\ne0$ preserves the mode but is \emph{not elastic};
only $m=n,q=0$ is elastic. A homogeneous drive supplies no compact
momentum and does not rotate one free particle into another mass branch.

\section{Readout with an unobserved supply channel}
In the driven prediction use
\begin{equation}
 h_{jq}(x)=w_{j+m-n}\int_0^\infty\dd p\,f_{\sigma_i}(p)
 K_{j+m-n,m,q}(p){\cal R}_{\sigma_d}(x\mid p_{f,jq}),\qquad
 h_j(x)=\sum_qh_{jq}(x).
\end{equation}
The classical beam is diagonal in $p$, and different $l$ give
orthogonal recoil states. With the long-time sideband limit this
justifies summing rates. It is not a treatment of finite coherent
wavepacket interference. If a hypothetical independent readout also
provided $q$, the optimal error would obey
\begin{equation}
 e_*^{\rm tagged}=\frac1{Z_D^{(m,r)}}\sum_q
 \int\dd x\,\min(h_{+q},h_{-q})\ \leq\ e_*^{\rm untagged}.
\end{equation}
This is an information bound, not a constructed pump measurement.
The passive inversion is invalid without accounting for $q$.
For known $p,q$, repair it by subtracting work:
\begin{equation}
 l^2=R^2[(E_f-q\Omega-E_n)^2-M_D^2-p^2],\qquad
 j=\frac{l^2-(m-n)^2-r^2}{2(m-n)} .
\end{equation}
At $p=.2$ the six momenta for the selected output are
\begin{center}
\begin{tabular}{rrrr}
\toprule
$j$ & $q=-1$ & $q=0$ & $q=+1$\\
\midrule
$+1$ & .82821583 & 1.02175458 & 1.20257369\\
$-1$ & 1.18924247 & 1.36191317 & 1.52787997\\
\bottomrule
\end{tabular}
\end{center}
The closest opposite-sign pair is separated by $.01333122$,
versus the passive carrier gap $.34015860$. All six are still
distinct: zero-noise catalogue identification at known $p$ remains
possible. At finite resolution use the full mixture, not just
the smallest spacing. Exact crossing of these two sidebands occurs at
\begin{equation}
 2\Omega_*=\sqrt{M_D^2+4/R^2+p^2}-\sqrt{M_D^2+p^2},
 \qquad \Omega_*(.2)=.19243951642 .
\end{equation}
This diagnostic is not used to tune the chosen drive. At the crossing,
those opposite-sign contributions have identical momentum readouts
unless additional information is measured.

\section{Completion boundary and the next physical decision}
The sharp-beam driven ledger contains 66 open $(l,m,q)$ rows, split
as $(14,25,27)$ across $q=(-1,0,+1)$. It gives
$Z_D=2.10563491267\,10^{-6}$, a mode-conversion fraction $.87381701$,
and mean pump work $.00900097215$ per counted scattering event.
These are not absolute encounter efficiencies or a power.

Selected optimal errors, in percent, illustrate the physical decision:
\begin{center}
\begin{tabular}{rrrrr}
\toprule
$\sigma_i$ & $\sigma_d$ & No drive & Drive, $q$ unknown & Drive, ideal $q$ tag\\
\midrule
0 & .1 & 3.55155 & 5.76199 & 3.55366\\
.1 & .1 & 3.77384 & 5.94485 & 3.77589\\
.2 & .1 & 4.84993 & 6.83823 & 4.85210\\
.4 & .1 & 11.0294 & 12.2710 & 11.0331\\
\bottomrule
\end{tabular}
\end{center}
Thus at $(.1,.1)$ the chosen 5\% criterion passes without a drive,
fails with unobserved pump exchange, and passes with an ideal tag.
This does not make every driven setup worse, or supply a physical tag.
At $\sigma_i=.1$, improving $\sigma_d$ to $.05$ instead gives a
driven untagged error of $2.20225\%$.

The reports retain all 50 resolution cards, conditional false-sign
errors, selected priors and decision boundaries. Quadrature orders
24 and 48 and two boundary-search mesh widths are compared. Gaussian
CDF intervals, not sampled histogram areas, give the error integrals.
For the retained incident region $z\leq10$, the omitted probability is
$\tau=\Phi_{\!N}(-10)/\Phi_{\!N}(p_0/\sigma_i)$. Since
$p_f/E_{\rm out}\leq1/2$,
\begin{equation}
 K_{lmq}\leq\frac{|g_q|^2}{32\pi m_n\mu_l},\qquad
 \delta Z\leq\tau\sum_{l\in\{0,-2\},q}
 \frac{w_l|g_q|^2}{32\pi m_n\mu_l}.
\end{equation}
An error-change bound is $2\delta Z/Z_{\rm retained}$; the
response-tail loss is at most $2\Phi_{\!N}(-10)$.
Mesh agreement is a numerical convergence check, not a theorem that
all possible mixture roots have been isolated for arbitrary cards.
Gaussian tails are
numerically bounded, not declared evidence that a five-dimensional
EFT is valid at arbitrarily high energies; a UV validity scale and
loop uncertainties remain unspecified. Results are tree-level,
long-time conditional forward predictions, not a physical particle fit.

The next bounded decision is an actual readout specification:
resolution, beam spread, mass confusion and any independent pump-energy
information. Absolute four-dimensional collision probabilities are
unnecessary for relative recoil discrimination. If finite pulses,
coherent interference or absolute probabilities are desired, specify
spatial envelopes, timing and relative pump phase together and compute
their common amplitude. A momentum histogram alone does not supply
those inputs.

\begin{thebibliography}{9}
\bibitem{pdg} Particle Data Group, \emph{Kinematics}, Review of Particle
Physics (2024), sections 49.3--49.5,
\url{https://pdg.lbl.gov/2024/reviews/rpp2024-rev-kinematics.pdf}.
\bibitem{floquet} M. Moskalets and M. B\"uttiker,
\emph{Floquet scattering theory of quantum pumps},
Phys. Rev. B \textbf{66}, 205320 (2002),
\url{https://arxiv.org/abs/cond-mat/0208356}.
\end{thebibliography}
\end{document}
